%% ============================================================
\section{Method}\label{sec:method}
%% ============================================================

CRAFT shares one clause-processing pipeline across inference, SFT target
construction, and RL credit assignment (Figure~\ref{fig:overview}).
Appendix~\ref{app:formal-method} gives the complete formal definition.

\subsection{Target-Hidden Composable Inference}
\label{sec:method:core}
\label{sec:method:inference}

Clauses from different responses can provide complementary target
constraints and support one another during inductive verification.
We therefore pool them before filtering.  Let
\(\operatorname{Admit}(y)\) extract the first \(m_{\max}\) invariant clauses
from response \(y\) and remove duplicates, with \(m_{\max}\) specified in
Table~\ref{tab:constants}.  Algorithm~\ref{alg:method-pipeline} gives the
inference procedure; \(\Vert\) denotes concatenation and
\(\operatorname{Unique}\) removes duplicates.

\begin{algorithm}[H]
\caption{Rollout--decompose--pool--filter--compose}
\label{alg:method-pipeline}
\begin{algorithmic}
\alginput{task \(\mathcal T=(\mathcal L,\mathcal G)\), policy \(\pi\), number of responses \(n\)}
\algoutput{composed invariant \(I\), final verdict \(v\)}
\algline \(x\leftarrow\mathcal L\) by masking \(\mathcal G\) in \(\mathcal T\)\cmnt{hide target}\\
\algline \(y_{1:n}\sim\pi(\cdot\mid x)\)\cmnt{rollout}\\
\algline \(A_i\leftarrow\operatorname{Admit}(y_i)\) for \(i=1,\ldots,n\)
  \cmnt{decompose and cap}\\
\algline \(P\leftarrow\operatorname{Unique}(A_1\Vert\cdots\Vert A_n)\)
  \cmnt{pool in response order}\\
\algline \(C\leftarrow\operatorname{Filter}_{\mathcal L}(P)\)\cmnt{filter}\\
\algline \(I\leftarrow\bigwedge_{c\in C}c\)\cmnt{compose}\\
\algline \(v\leftarrow\operatorname{TaskVerify}(\mathcal T,C)\)\cmnt{restore target}\\
\algline \kw{return} \((I,v)\)\\
\end{algorithmic}
\end{algorithm}

\tightparagraph{Joint verification and deletion.}
Houdini checks each clause for establishment at loop entry and preservation
under the current conjunction.  It repeatedly deletes unproved clauses and
rechecks the survivors, since they may have depended on removed clauses.
Only a proved fixed point is returned, ensuring an inductive conjunction
under a sound verifier.  Budget and failure handling are specified in
Appendix~\ref{app:formal-houdini}.

\tightparagraph{Why pooling precedes filtering.}
Let \(\mathsf H_{\mathcal L}\) denote ideal Houdini with exact logical
decisions.  Filtering the union retains every clause obtainable by separately
filtering the responses:
\begin{equation}
 \bigcup_i\mathsf H_{\mathcal L}(A_i)
 \subseteq\mathsf H_{\mathcal L}\!\left(\bigcup_i A_i\right).
 \label{eq:main-pooling-inclusion}
\end{equation}
The inclusion can be strict when supporting clauses occur in different
responses; its derivation is in Appendix~\ref{app:composition-proofs}.
Finite prover budgets can break this inclusion, so measured compose@\(k\)
need not be monotonic (Appendix~\ref{app:houdini-order}).

\subsection{Target-Independent Strength Signal}
\label{sec:method:sampler}

Inductive validity alone cannot distinguish useful invariants from weak
ones such as \(\mathsf{true}\).  We execute the masked loop and construct
candidate negative traces \(\mathcal N\): local relational perturbations
probe constraints among variables, while continuations beyond witnessed
exits probe boundary constraints.  Sampling and cleaning are detailed in
Appendix~\ref{app:negative-sampling}.

For a verified clause set \(C\), the state evaluator returns
\(\operatorname{Eval}(c,s)\in\{\mathsf{true},\mathsf{false},\mathsf{unknown}\}\).
A negative trace is rejected only by a definite false evaluation:
\begin{equation}
 \operatorname{rej}(C)=\{\nu\in\mathcal N:\exists s\in\nu,\exists c\in C,
   \operatorname{Eval}(c,s)=\mathsf{false}\}.
 \label{eq:main-negative-rejection}
\end{equation}
An unknown evaluation contributes no rejection.  Perturbations form
singleton traces; each cleaned post-exit continuation remains one trace.
The trajectory-level coverage is
\begin{equation}
  \operatorname{cov}(C)=
  \frac{|\operatorname{rej}(C)|}{|\mathcal N|},
  \qquad |\mathcal N|>0.
  \label{eq:method-coverage}
\end{equation}
Each trace contributes one unit regardless of its length.

\tightparagraph{Empirical alignment.}
The verifier supplies soundness, while coverage provides an empirical strength
signal. Final verification determines hidden-target sufficiency. Coverage
achieves a within-program macro
AUROC of 0.842 for final verification, supporting this association;
Appendix~\ref{app:coverage-predictiveness} gives the diagnostic population,
uncertainty, and limitations.

\subsection{Distilling Multi-Response Search}
\label{sec:method:sft}

SFT distills accepted multi-response compositions into single-response
supervision, aiming to make a richer set of useful clauses available with
fewer samples.  Starting with \(P_0=\varnothing\), each round samples \(g_{\mathrm{sft}}\) responses
from the same base model that is subsequently fine-tuned, adds their clauses
to the accumulated pool, and jointly filters it:
\begin{equation}
  P_t=\operatorname{Unique}\!\left(P_{t-1}\Vert
    \operatorname{Admit}(y_1^{(t)})\Vert\cdots\Vert
    \operatorname{Admit}(y_{g_{\mathrm{sft}}}^{(t)})\right),
  \qquad
  C_t=\operatorname{Filter}_{\mathcal L}(P_t).
  \label{eq:method-sft}
\end{equation}
Accumulating all admitted clauses in \(P_t\) allows a removed clause to return
when a later response supplies support.
The target can therefore combine discoveries across responses and rounds.

\tightparagraph{Acceptance and serialization.}
For nonempty \(\mathcal N\), the first composition satisfying
\(\operatorname{cov}(C_t)\ge\tau\) within \(t_{\mathrm{sft}}\) rounds is
serialized as the SFT target; otherwise no training pair is produced.
Standard autoregressive cross-entropy trains the model on this target,
conditioned on the masked loop.  Algorithm~\ref{alg:formal-sft} gives the
complete procedure.

\subsection{Rewarding Decomposed Contributions}
\label{sec:method:reward}

A failed response should not erase the contributions of its useful clauses.
RL therefore assigns credit to the parts retained by joint filtering.
For a group sampled from the current policy, let
\(A_i=\operatorname{Admit}(y_i)\).  We pool and filter once, then map
surviving clauses back to their proposing responses:
\begin{align}
  P^*&=\operatorname{Unique}\!\left(A_1\Vert\cdots\Vert A_{g_{\mathrm{rl}}}\right),
  &C^*&=\operatorname{Filter}_{\mathcal L}(P^*),\\
  V_i&=\{c\in A_i:\operatorname{key}(c)\in\operatorname{keys}(C^*)\},
  &R_i&=\operatorname{rej}(V_i).
  \label{eq:method-credit-pool}
\end{align}
Canonical keys identify which clauses each response proposed.  Thus \(V_i\)
retains its surviving contributions even when other clauses fail, with
mutually supporting clauses checked in the pooled context used at inference.

For nonempty \(\mathcal N\), Full, our default reward, adds Shapley credit
to Clause-decomposed coverage:
\begin{equation}
\label{eq:generation-reward}
  r_i=\frac{|R_i|}{|\mathcal N|}+0.3\phi_i.
\end{equation}
With the survivor bundles fixed, let \(f(\nu)=|\{j:\nu\in R_j\}|\) count
the responses rejecting trace \(\nu\).  The Shapley credit is
\begin{equation}
  \phi_i=\frac{1}{|\mathcal N|}\sum_{\nu\in R_i}\frac{1}{f(\nu)}.
  \label{eq:method-shapley}
\end{equation}
Each rejected trace's credit is divided among the responses covering it.
At equal coverage, Full rewards responses more when their rejected traces
are covered by fewer other responses, encouraging complementary effective
coverage within the group.
The coverage game and proof are in Appendix~\ref{app:formal-shapley};
Appendix~\ref{app:full-credit-analysis} analyzes the empirical effect.

GRPO standardizes rewards within each group and applies each response's
scalar advantage to its tokens; clauses determine credit, while policy
updates remain response-level.  A clause that only supports another clause
receives no direct coverage credit.  Update details and the empty-negative
fallback are in Appendix~\ref{app:grpo-details} and
Algorithm~\ref{alg:formal-rl-reward}.
